\section*{Capítulo \ref{ch:b3:plant-molecular-physiology} --- Fisiologia vegetal molecular e respostas ao estresse}

\begin{solution}{exo:b3:plant-molecular-physiology:1}
\omterm{def:b3:plant-molecular-physiology:photoreceptors}{Fitocromos}, vermelho \qty{660}{nm} e vermelho-distante \qty{730}{nm}:
desestiolamento, \omterm{prop:b3:plant-molecular-physiology:photoequilibrium}{fuga da sombra}, germinação. \omterm{def:b3:plant-molecular-physiology:photoreceptors}{Criptocromos}, azul
\qty{450}{nm}: desestiolamento, relógio, floração. \omterm{def:b3:plant-molecular-physiology:photoreceptors}{Fototropinas}, azul:
curvatura em direção à luz, abertura estomática, movimento dos
cloroplastos. A clorofila mede a luz como energia e nada faz com a
informação; os fotorreceptores leem cor, direção e duração em fluxos um
milhão de vezes menores e mudam a expressão gênica --- a diferença entre
comer a luz e lê-la.
\end{solution}

\begin{solution}{exo:b3:plant-molecular-physiology:2}
Cada \omterm{def:b3:endocrinology:hormone}{hormônio} promove a ubiquitinação e a destruição de um repressor: a
\omterm{def:b3:plant-molecular-physiology:hormones}{auxina} cola a \omterm{def:b3:plant-molecular-physiology:hormones}{TIR1} às proteínas Aux/IAA, a \omterm{def:b3:plant-molecular-physiology:hormones}{giberelina} ligada a GID1
entrega as proteínas \omterm{def:b3:plant-molecular-physiology:hormones}{DELLA} a uma ligase F-box, o jasmonato cola a COI1
às proteínas JAZ. Os repressores já estão sobre seus genes-alvo, com os
ativadores prontos ao lado, de modo que nenhuma proteína nova precisa ser
feita: destruir o repressor, o que o proteassomo faz em minutos, liga os
genes de imediato.
\end{solution}

\begin{solution}{exo:b3:plant-molecular-physiology:3}
O ABA liga PYR/PYL; o complexo inibe a fosfatase PP2C; a quinase SnRK2
(OST1), já não desfosforilada, fica ativa e fosforila o canal de ânions
SLAC1; os ânions saem, a membrana despolariza, os canais de efluxo de
potássio abrem, o K$^{+}$ e depois a água saem, a turgescência cai e as
células-guarda murcham uma contra a outra. Qualquer elo pode ser rompido:
sem receptores, uma fosfatase que não pode ser inibida (o clássico
mutante \emph{abi1}), sem OST1 ou sem SLAC1 --- cada um dá uma planta
murcha cujos estômatos ficam abertos na seca.
\end{solution}

\begin{solution}{exo:b3:plant-molecular-physiology:4}
A PTI reconhece moléculas comuns a classes microbianas inteiras
(flagelina, quitina) com quinases receptoras na superfície celular e
termina numa parede reforçada, estômatos fechados, expressão de genes de
defesa e crescimento mais lento, em geral sem morte celular. A ETI
reconhece uma proteína efetora específica ou o dano que ela faz, com um
\omterm{def:b3:plant-molecular-physiology:immunity}{receptor NLR} intracelular, e termina na morte hipersensível das células
infectadas e num alarme sistêmico.
\end{solution}

\begin{solution}{exo:b3:plant-molecular-physiology:5}
Luz do dia, $0.87\times 1.15/1.75 = 0.57$; uma folha, $0.87\times 0.4/1.0 =
0.35$; dossel denso, $0.87\times 0.1/0.7 = 0.12$; lâmpada
incandescente, $0.87\times 0.7/1.3 = 0.47$. Uma lâmpada de filamento é
rica em vermelho-distante, de modo que $\varphi$ fica abaixo do valor da
luz do dia e a plântula lê sombra parcial: ela se alonga. As lâmpadas
brancas frias e os LEDs sem vermelho-distante fazem o oposto.
\end{solution}

\begin{solution}{exo:b3:plant-molecular-physiology:6}
Parede $5.0$, citosol $7.5$: $(1 + 10^{2.75})/(1 + 10^{0.25}) = 563/2.78
= 203$. Parede $4.5$: $563/(1 + 10^{-0.25}) = 563/1.56 = 360$ ---
acidificar a parede eleva a fração neutra do lado de fora e a captação.
Citosol $6.5$ com parede $5.5$: $(1 + 10^{1.75})/6.62 = 57/6.62 = 8.6$: o
aprisionamento cai cinco vezes, a \omterm{def:b3:plant-molecular-physiology:hormones}{auxina} vaza por todas as faces e os
gradientes polares se achatam --- uma razão pela qual raízes em anoxia
perdem a orientação.
\end{solution}

\begin{solution}{exo:b3:plant-molecular-physiology:7}
\qty{1}{cm/h} $= \qty{10000}{\micro m/h}$: $200$ células por hora,
\qty{18}{s} em cada uma. Difusão através de \qty{50}{\micro m}: $t =
(5\times 10^{-5})^{2}/(2\times 5\times 10^{-10}) = \qty{2.5}{s}$. O
interior da célula é atravessado em segundos; a maior parte dos
\qty{18}{s} é passada esperando para ser carregada para fora pelas PIN
da membrana basal, que é a etapa limitante.
\end{solution}

\begin{solution}{exo:b3:plant-molecular-physiology:8}
$E/A = 1.6\,\Delta w/\Delta c = 1.6\times\num{16000}/150 = 171$ moléculas
de água por CO$_{2}$. C$_{4}$, $\Delta c = 250$: $102$. Dia seco,
$\Delta w = 24$: $256$.
\end{solution}

\begin{solution}{exo:b3:plant-molecular-physiology:9}
O congelamento desidrata as células ao puxar água para o gelo
extracelular, de modo que o frio e a seca são ambos estresses de
desidratação. Eles compartilham segundos mensageiros (ABA, picos de
cálcio, oxigênio reativo), fatores de transcrição (a família CBF/DREB é
induzida pelos dois) e produtos: desidrinas e proteínas LEA que protegem
membranas e proteínas da perda de água, solutos compatíveis (prolina,
açúcares) que retêm água, e antioxidantes. Uma planta que os fez para um
dos estresses já os tem para o outro.
\end{solution}

\begin{solution}{exo:b3:plant-molecular-physiology:10}
Em \qty{1}{\%} do sol pleno, $k_{1} + k_{2} = \qty{0.005}{s^{-1}}$:
95~\% do equilíbrio depois de $3/k = \qty{600}{s}$, dez minutos, contra
seis segundos ao meio-dia. Oito horas são quatro meias-vidas: resta
$1/16 \approx
\qty{6}{\%}$ da Pfr. Um lampejo breve fixa $\varphi$, mas a
Pfr então decai, ao passo que um dia a mantém alta por horas; as
respostas que exigem que a Pfr aja por horas (germinação,
desestiolamento) integram a exposição, de modo que uma semente revirada
por um arado num instante não responde como uma que jaz na superfície.
\end{solution}

\begin{solution}{exo:b3:plant-molecular-physiology:11}
Dez sítios por dia a $100$ células cada: \num{1000} células, $10^{-4}$ da
folha; ao longo de uma estação, uma fração de um por cento, contra a
perda da folha se uma infecção se espalhasse sem controle. As células
mortas custam quase nada e levam consigo o alimento do patógeno --- para
um biotrófico, que precisa de células vivas. Um necrotrófico se alimenta
de tecido morto: a \omterm{def:b3:plant-molecular-physiology:immunity}{resposta hipersensível} lhe entrega uma refeição e um
ponto de apoio, e alguns (\emph{Botrytis}, \emph{Sclerotinia}) secretam
toxinas que a provocam de propósito; contra eles a planta conta com as
defesas mediadas por jasmonato, e não com a morte celular.
\end{solution}

\begin{solution}{exo:b3:plant-molecular-physiology:12}
Produção à taxa $p$ da hora 12 à 16. Dia longo, com luz o tempo todo,
$k = \ln 2/4 = \qty{0.17}{h^{-1}}$: o nível na hora 16 é $(p/k)(1 -
\mathrm{e}^{-4k}) = (p/0.17)(0.5) = 2.9p$. Dia curto, com escuro a partir
da hora 10, $k = \qty{1.39}{h^{-1}}$: $(p/1.39)(1 - \mathrm{e}^{-5.5})
\approx 0.72p$. Quatro vezes mais CONSTANS no dia longo. Um limiar
intermediário --- digamos $2p$, necessário para ativar o \emph{FT} --- só
é cruzado quando a janela de produção do relógio coincide com a luz: a
coincidência de um ritmo interno com o dia externo converte uma duração
graduada de dia numa decisão tudo ou nada de florescer.
\end{solution}

\begin{solution}{pb:b3:plant-molecular-physiology:1}
\textbf{1.} Luz do dia, $0.87\times 1.15/1.75 = 0.57$; dossel,
$0.87\times
0.2/0.8 = 0.22$.
\textbf{2.} $r = 2(1 - \varphi)$: \qty{0.86}{mm/h} à luz do dia,
\qty{1.57}{mm/h} na sombra; em \qty{48}{h}: $41$ contra \qty{75}{mm},
\qty{34}{mm} a mais.
\textbf{3.} O vermelho cai a $0.1\times 1.15 = 0.115$ do
vermelho-distante, que cai a $0.8$: $\zeta = 0.144$, $\varphi = 0.87\times 0.144/0.744 =
0.17$. Sim: uma única folha derruba $\varphi$ de
$0.57$ para $0.17$, bem dentro da faixa de \omterm{prop:b3:plant-molecular-physiology:photoequilibrium}{fuga da sombra}.
\textbf{4.} As duas taxas de fotoconversão escalam com o fluxo, de modo
que sua razão não escala. Sob nuvens o espectro quase não muda e
$\varphi$ fica perto de $0.57$: uma plântula a céu aberto num dia
nublado não se alonga --- ela lê vizinhos, e não brilho.
\textbf{5.} \qty{95}{\%} depois de $3/(k_{1} + k_{2})$: \qty{6}{s} no sol
pleno, \qty{600}{s} em \qty{1}{\%}.
\textbf{6.} \qty{8}{h} $=$ quatro meias-vidas, $1/16 = \qty{6.3}{\%}$;
\qty{14}{h} $=$ sete, $1/128 = \qty{0.8}{\%}$. A Pfr residual ao
amanhecer poderia informar a duração da noite, mas a reversão é uma
reação térmica que corre mais depressa em noites quentes, e o valor
inicial depende da luz do dia: as plantas medem a duração da noite com o
relógio, em vez disso.
\textbf{7.} Parede: $1/(1 + 10^{0.75}) = 0.15$; citosol: $1/(1 +
10^{2.45}) = 0.0035$.
\textbf{8.} $(1 + 282)/(1 + 5.6) = 283/6.6 = 43$. Citosol,
$43\times 0.1 = \qty{4.3}{\micro M}$.
\textbf{9.} $\qty{1}{cm/h}/\qty{100}{\micro m} = 100$ células por hora,
\qty{36}{s} em cada.
\textbf{10.} \qty{5}{h}. A curvatura começa em menos de uma hora porque a
redistribuição que importa é lateral, ao longo de um milímetro de ponta,
e a resposta é local; o fluxo de longa distância fixa o suprimento, e não
o \omterm{thm:b3:rna-regulation:mirna-kinetics}{tempo de resposta}.
\textbf{11.} Flanco inferior: $60/50 = 1.2$, \omterm{def:b3:plant-molecular-physiology:hormones}{auxina} \qty{20}{\%} acima da
norma, alongamento $0.8$; flanco superior: $0.8$ da norma, alongamento
$1.2$. O flanco superior cresce mais que o inferior: a raiz se curva para
baixo.
\textbf{12.} As células do caule estão abaixo de seu ótimo de \omterm{def:b3:plant-molecular-physiology:hormones}{auxina}, de
modo que a \omterm{def:b3:plant-molecular-physiology:hormones}{auxina} extra do flanco inferior promove seu crescimento: o
flanco inferior cresce mais que o superior e o caule se curva para cima.
\textbf{13.} $E = 0.3\times 16 = \qty{4.8}{mmol\,m^{-2}\,s^{-1}}$.
\textbf{14.} $A = (0.3/1.6)\times 150 = \qty{28}{\micro mol\,m^{-2}\,s^{-1}}$;
$E/A = 4800/28 = 171$.
\textbf{15.} Área \qty{0.002}{m^2}, \num{43200}~s: água $4.8\times
10^{-3}\times 0.002\times\num{43200} = \qty{0.41}{mol} = \qty{7.5}{g}$;
CO$_{2}$ $28\times 10^{-6}\times 0.002\times\num{43200} =
\qty{2.4e-3}{mol}$, isto é, $2.4\times 10^{-3}/6\times 180 =
\qty{0.073}{g}$ de glicose --- cem gramas de água por grama de açúcar.
\textbf{16.} $E = 0.05\times 16 = \qty{0.8}{mmol\,m^{-2}\,s^{-1}}$,
$A = \qty{4.7}{\micro mol\,m^{-2}\,s^{-1}}$, razão ainda $171$. A planta
poupou seis sétimos de sua água e abriu mão de seis sétimos de seu
carbono: fechar a válvula muda a taxa, e não o preço.
\textbf{17.} C$_{4}$: $\Delta c = 250$, $E/A = 4800/46.9 = 102$. CAM à
noite: $E = 0.3\times 5 = \qty{1.5}{mmol}$, $A = 28$: $E/A = 53$.
\textbf{18.} $E/A = 1.6\,\Delta w/\Delta c$: a condutância se cancela
porque os dois gases passam pelo mesmo poro. A planta pode mudar os
gradientes --- concentrar CO$_{2}$ por dentro (C$_{4}$), abrir quando o
ar está úmido (CAM, amanhecer), resfriar a folha, engrossar a camada
limite com pelos ---, mas não a física de dois gases que dividem um
buraco.
\textbf{19.} \num{1000} células por dia, $10^{-4}$ da folha.
\textbf{20.} \num{60000} células em 60 dias, \qty{0.6}{\%}. Uma escapada
por dia destruindo \qty{5}{\%} consumiria a folha em vinte dias: a
\omterm{def:b3:plant-molecular-physiology:immunity}{resposta hipersensível} custa um centésimo do que custa uma única
infecção que escapa.
\textbf{21.} $0.2\times 5 = \qty{1}{\%}$ da fotossíntese. Mantida sempre
ligada custaria \qty{5}{\%} e a planta seria ultrapassada por vizinhas
que gastassem isso em folhas; a defesa induzida só compensa quando a
ameaça está presente.
\textbf{22.} Ele ganha invisibilidade para aquele NLR e infecta a
variedade resistente; perde o que quer que o efetor fizesse (suprimir a
PTI), de modo que fica mais fraco em plantas sem o NLR. A população de
plantas passa a favorecer NLRs contra os efetores restantes, e o velho
gene de resistência, agora inútil, declina: uma ciclagem dependente da
frequência dos genes dos dois lados.
\textbf{23.} O imitador (coronatina) ativa a via do jasmonato, que
antagoniza a via do salicilato; o salicilato é a defesa contra
biotróficos como a própria bactéria, e o imitador ainda reabre estômatos
que a planta havia fechado. O patógeno vira um braço do sistema imune
contra o outro.
\textbf{24.} Um único gene de resistência em milhões de plantas
idênticas é uma seleção uniforme: qualquer mutante que perca ou altere o
efetor se espalha por toda a lavoura em poucas safras. Os melhoristas
empilham vários genes de resistência, de modo que vários efetores tenham
de ser perdidos ao mesmo tempo, misturam variedades, rodiziam genes ao
longo dos anos e favorecem a ampla resistência desencadeada por padrões,
de que nenhuma mutação isolada escapa.
\textbf{25.} $\varphi \approx 0.22$ sob o dossel; \omterm{def:b3:plant-molecular-physiology:hormones}{auxina} dentro :
fora $\approx 43$; cerca de $170$ moléculas de água perdidas por
CO$_{2}$ fixado.
\end{solution}
